% ---
% titre: Représenter une fraction
% theme: NO
% chapitre: Nombres rationnels
% sous_chapitre: Représenter, comparer et encadrer
% niveau: [10e]
% type: EVAL
% tags: [c-fractions, c-representer, p-question-TS]
% difficulte: 1
% corrige: true
% ---

\newcommand{\GrilleFig}[6]{%
\begin{tikzpicture}[baseline=(current bounding box.center)]
\pgfmathsetmacro{\gfw}{#3/#1}
\pgfmathsetmacro{\gfh}{#4/#2}
\foreach \gfr in {0,...,\numexpr#2-1\relax}{
  \foreach \gfc in {0,...,\numexpr#1-1\relax}{
    \pgfmathtruncatemacro{\gfi}{\gfr*#1+\gfc}
    \ifnum\gfi<#5\relax\def\gffill{#6}\else\def\gffill{white}\fi
    \draw[fill=\gffill] (\gfc*\gfw,-\gfr*\gfh) rectangle ++(\gfw,-\gfh);
  }}
\end{tikzpicture}}

\newcommand{\DisqueFig}[3]{%
\begin{tikzpicture}[baseline=(current bounding box.center)]
\foreach \dfs in {0,...,\numexpr#1-1\relax}{
  \pgfmathsetmacro{\dfa}{90-360/#1*\dfs}
  \pgfmathsetmacro{\dfb}{90-360/#1*(\dfs+1)}
  \ifnum\dfs<#2\relax\def\dffill{#3}\else\def\dffill{white}\fi
  \draw[fill=\dffill] (0,0) -- (\dfa:1) arc[start angle=\dfa, end angle=\dfb, radius=1] -- cycle;
}
\end{tikzpicture}}


\question
\begin{parts}
\part[2] Colorie la fraction indiquée de chaque figure.

\vspace{10pt}
\def\arraystretch{1.5}
\begin{tabularx}{\linewidth}{XX}
$\dfrac{3}{5}$ \quad \GrilleFig{5}{1}{4.5}{1.2}{0}{white} &$\dfrac{5}{8}$ \quad \GrilleFig{4}{2}{3.6}{1.2}{0}{white} \\[26pt]
$\dfrac{2}{3}$ \quad \DisqueFig{3}{0}{white} &$\dfrac{3}{8}$ \quad \DisqueFig{8}{0}{white} \\
\end{tabularx}
\def\arraystretch{1}

\vspace{15pt}
\part[2]  Écris la fraction correspondant à la partie coloriée.

\vspace{10pt}
\def\arraystretch{1.5}
\begin{tabularx}{\linewidth}{XcX}
\GrilleFig{5}{1}{4.5}{1.2}{4}{themeColor!40} \quad $\Longrightarrow$ \sol{\;\bl}{$\dfrac{4}{5}$} && \DisqueFig{4}{3}{themeColor!40} \quad $\Longrightarrow$ \sol{\;\bl}{$\dfrac{3}{4}$} \\[26pt]
\GrilleFig{3}{2}{3.6}{1.6}{4}{themeColor!40} \quad $\Longrightarrow$ \sol{\;\bl}{$\dfrac{4}{6}$} && \GrilleFig{5}{2}{4.5}{1.6}{7}{themeColor!40} \quad $\Longrightarrow$ \sol{\;\bl}{$\dfrac{7}{10}$} \\
\end{tabularx}
\def\arraystretch{1}

\begin{solution}[0pt]
\textbf{0.5pt par réponse}\\[4pt]
\textcolor{themeColor}{Partie A.} a) $3$ cases sur $5$ \quad b) $5$ cases sur $8$ \quad c) $2$ parts sur $3$ \quad d) $3$ parts sur $8$
\end{solution}
\end{parts}